\section{问题定义与业务时间对齐}

设预测日为 \(D\)，数据时间分辨率为 1 小时。为区分气象产品的起报时刻和功率预测起点，记 \(t_{\mathrm{init}}\) 为 Pangu-Weather 的气象起报时刻，即北京时间 \(D-1\) 日 08:00；记 \(t_0\) 为功率预测起点，即 \(D-1\) 日 23:00。本文以 \(L=16\) 表示功率状态窗口包含的逐小时观测数，该窗口覆盖 \(t_0-L+1\) 至 \(t_0\)，即 \(D-1\) 日 08:00--23:00；以 \(H=24\) 表示日内预测步数，第 \(h\) 个预测步对应 \(t_0+h\)，\(h=1,\ldots,H\)，因而目标窗口覆盖 \(D\) 日 00:00--23:00。

记 \(N\) 为区域内风电场站数量，预测起点前的状态上下文为 \(\mathbf{P}^{c}\in\mathbb{R}^{N\times L}\)。区域聚合功率目标记为 \(\mathbf{y}=[y_{t_0+1},\ldots,y_{t_0+H}]\in\mathbb{R}^{H}\)，模型输出为 \(\hat{\mathbf{y}}\in\mathbb{R}^{H}\)。场站邻接矩阵 \(\mathbf{A}\in\mathbb{R}^{N\times N}\) 由同一 \(L\) 步功率状态窗口估计，保留 Pearson 相关系数大于 0.7 的正相关系数作为边权。

业务时间协议如图~\ref{fig:time-protocol} 所示。由 \(t_{\mathrm{init}}\) 起报的同一次 Pangu-Weather 产品中，预报时效 0 至 \(L-1\) 对应 \(D-1\) 日 08:00--23:00，记为气象预报历史段 \(\mathbf{X}^{h}\)；预报时效 \(L\) 至 \(L+H-1\) 对应目标窗口 \(D\) 日 00:00--23:00，记为气象预报目标段 \(\mathbf{X}^{f}\)。这里的“历史段”和“目标段”仅表示气象帧的有效时刻位于目标窗口之前或之内，二者均来自 \(t_{\mathrm{init}}\) 时已经生成的同一次气象预报。两段沿时间维拼接为 \(\mathbf{X}=[\mathbf{X}^{h};\mathbf{X}^{f}]\)。因此，以 \(t_0\) 为起点生成的功率预测只读取截至 \(t_0\) 的实测功率和此前起报的气象产品，不使用目标窗口内的实测功率。

训练阶段采用逐小时滑动窗口，并保持上述输入窗口与目标窗口的相对时间对齐；验证和测试阶段采用严格日样本，按照固定业务时间协议预测 \(D\) 日出力。
